\section{Clubs and stationary subsets of ordinals}

\import{cardinal/normal-map.tex}
\import{set/indexed-family.tex}
\import{set/iteration.tex}

All predicates in this section carry their ambient ordinal explicitly.  In
particular, a club is not a global class object and stationarity is a
predicate on a bounded subset. Diagonal intersection and normal-map preimage
are deliberately total bounded trace constructors; their mathematical laws
provide the ordinal and function hypotheses where needed.

\begin{definition}\label{ordinal_unbounded}
    $\ordinalUnbounded{C}{\alpha}$ iff $C\subseteq\alpha$ and
        for every $\beta\in\alpha$ there exists $\gamma\in C$ such that
            $\beta\in\gamma$.
\end{definition}

\begin{definition}\label{ordinal_closed}
    $\ordinalClosed{C}{\alpha}$ iff $C\subseteq\alpha$ and
        for every limit ordinal $\gamma\in\alpha$ such that
            $\ordinalUnbounded{C\inter\gamma}{\gamma}$
        we have $\gamma\in C$.
\end{definition}

\begin{definition}\label{club_in}
    $\clubIn{C}{\alpha}$ iff $\alpha$ is a limit ordinal and
        $\ordinalUnbounded{C}{\alpha}$ and $\ordinalClosed{C}{\alpha}$.
\end{definition}

\begin{definition}\label{stationary_in}
    $\stationaryIn{S}{\alpha}$ iff $\alpha$ is a limit ordinal and
        $S\subseteq\alpha$ and for every $C$ such that $\clubIn{C}{\alpha}$
            we have $S\meets C$.
\end{definition}

\begin{abbreviation}\label{stationary_in_phrase}
    $S$ is stationary in $\alpha$ iff $\stationaryIn{S}{\alpha}$.
\end{abbreviation}

\begin{definition}\label{nonstationary_in}
    $\nonstationaryIn{S}{\alpha}$ iff $\alpha$ is a limit ordinal and
        $S\subseteq\alpha$ and there exists $C$ such that
            $\clubIn{C}{\alpha}$ and $S\notmeets C$.
\end{definition}

\begin{definition}\label{ordinal_bounded}
    $\ordinalBounded{B}{\alpha}$ iff $B\subseteq\alpha$ and
        there exists $\beta\in\alpha$ such that $B\subseteq\beta$.
\end{definition}

\begin{definition}\label{ordinal_tail}
    $\ordinalTail{\beta}{\alpha}=\alpha\setminus\beta$.
\end{definition}

\begin{definition}\label{limit_ordinals_below}
    $\limitOrdinalsBelow{\alpha}
        =\{\gamma\in\alpha\mid \text{$\gamma$ is a limit ordinal}\}$.
\end{definition}

\begin{definition}\label{diagonal_intersection}
    $\diagonalIntersection{F}{\alpha}
        =\{\beta\in\alpha\mid
            \text{for every $i\in\beta$ we have $\beta\in F(i)$}\}$.
\end{definition}

\begin{definition}\label{normal_map_preimage}
    $\normalMapPreimage{c}{S}{\lambda}
        =\{i\in\lambda\mid c(i)\in S\}$.
\end{definition}

\begin{definition}\label{club_upper}
    $\clubUpper{C}{\beta}{\alpha}
        =\{\gamma\in C\mid \beta\in\gamma\}$.
\end{definition}

\begin{definition}\label{least_member_above}
    $\leastMemberAbove{C}{\beta}{\alpha}
        =\inters{\clubUpper{C}{\beta}{\alpha}}$.
\end{definition}

\begin{definition}\label{club_intermediate}
    $\clubIntermediate{C}{\beta}{\alpha}
        =\choice{\clubUpper{C}{\beta}{\alpha}}$.
\end{definition}

\begin{definition}\label{club_intersection_advance}
    $\clubIntersectionAdvance{C}{D}{\alpha}
        =\{(\beta,
            \choice{\clubUpper{D}{\clubIntermediate{C}{\beta}{\alpha}}{\alpha}})
            \mid \beta\in\alpha\}$.
\end{definition}

\begin{proposition}\label{club_in_elim}
    Suppose $\clubIn{C}{\alpha}$.
    Then $\alpha$ is a limit ordinal and $C\subseteq\alpha$ and
        $\ordinalUnbounded{C}{\alpha}$ and $\ordinalClosed{C}{\alpha}$.
\end{proposition}
\begin{proof}
    Follows by \auto{club_in,ordinal_unbounded,ordinal_closed}.
\end{proof}

\begin{proposition}\label{stationary_in_elim}
    Suppose $S$ is stationary in $\alpha$.
    Then $\alpha$ is a limit ordinal and $S\subseteq\alpha$.
\end{proposition}
\begin{proof}
    Follows by \obv{stationary_in}.
\end{proof}

\begin{proposition}\label{cofinal_map_into_limit_has_unbounded_range}
    Let $\alpha$ be a limit ordinal.
    Suppose $\cofinalMap{f}{D}{\alpha}$.
    Then $\ordinalUnbounded{\ran{f}}{\alpha}$.
\end{proposition}
\begin{proof}
    $f\in\funs{D}{\alpha}$ and $\cofinalIn{\ran{f}}{\alpha}$
        by \obv{cofinal_map}.
    Hence $\ran{f}\subseteq\alpha$ by \obv{funs_ran}.
    Show that for every $\beta\in\alpha$ there exists
        $\gamma\in\ran{f}$ such that $\beta\in\gamma$.
    \begin{subproof}
        Fix $\beta\in\alpha$.
        $\suc{\beta}\in\alpha$ by \obv{suc_elem_limitordinal}.
        Take $\gamma\in\ran{f}$ such that $\suc{\beta}\subseteq\gamma$
            by \obv{cofinal_in}.
        $\beta\in\suc{\beta}$ by \obv{suc_intro_self}.
        Hence $\beta\in\gamma$ by \obv{subseteq}.
    \end{subproof}
    Follows by \obv{ordinal_unbounded}.
\end{proof}

\begin{proposition}\label{ordinal_unbounded_iff_union}
    Let $\alpha$ be a limit ordinal.
    Suppose $C\subseteq\alpha$.
    Then $\ordinalUnbounded{C}{\alpha}$ iff $\unions{C}=\alpha$.
\end{proposition}
\begin{proof}
    Show that if $\ordinalUnbounded{C}{\alpha}$, then
        $\unions{C}=\alpha$.
    \begin{subproof}
        Assume $\ordinalUnbounded{C}{\alpha}$.
        Show that every element of $\unions{C}$ is an element of $\alpha$.
        \begin{subproof}
            Fix $x\in\unions{C}$.
            Take $\gamma\in C$ such that $x\in\gamma$
                by \obv{unions_iff}.
            Hence $\gamma\in\alpha$ by \obv{subseteq}.
            Follows by \obv{ordinal_transitivity,limit_ordinal}.
        \end{subproof}
        Hence $\unions{C}\subseteq\alpha$ by \obv{subseteq}.
        Show that every element of $\alpha$ is an element of $\unions{C}$.
        \begin{subproof}
            Fix $\beta\in\alpha$.
            Take $\gamma\in C$ such that $\beta\in\gamma$
                by \obv{ordinal_unbounded}.
            Follows by \obv{unions_iff}.
        \end{subproof}
        Hence $\alpha\subseteq\unions{C}$ by \obv{subseteq}.
        Follows by \obv{subseteq_antisymmetric}.
    \end{subproof}
    Show that if $\unions{C}=\alpha$, then
        $\ordinalUnbounded{C}{\alpha}$.
    \begin{subproof}
        Assume $\unions{C}=\alpha$.
        Show that for every $\beta\in\alpha$ there exists $\gamma\in C$
            such that $\beta\in\gamma$.
        \begin{subproof}
            Fix $\beta\in\alpha$.
            $\suc{\beta}\in\alpha$ by \obv{suc_elem_limitordinal}.
            Hence $\suc{\beta}\in\unions{C}$.
            Take $\gamma\in C$ such that $\suc{\beta}\in\gamma$
                by \obv{unions_iff}.
            $\beta\in\suc{\beta}$ by \obv{suc_intro_self}.
            $\gamma$ is an ordinal
                by \obv{subseteq,elem_of_ordinal_is_ordinal,limit_ordinal}.
            Hence $\beta\in\gamma$ by \obv{ordinal_transitivity}.
        \end{subproof}
        Follows by \obv{ordinal_unbounded}.
    \end{subproof}
    Follows by \auto.
\end{proof}

\begin{proposition}\label{ordinal_unbounded_superset}
    Suppose $\ordinalUnbounded{C}{\alpha}$.
    Suppose $C\subseteq D\subseteq\alpha$.
    Then $\ordinalUnbounded{D}{\alpha}$.
\end{proposition}
\begin{proof}
    For every $\beta\in\alpha$ there exists $\gamma\in C$ such that
        $\beta\in\gamma$ by \obv{ordinal_unbounded}.
    Hence for every $\beta\in\alpha$ there exists $\gamma\in D$ such that
        $\beta\in\gamma$ by \auto{subseteq}.
    Follows by \obv{ordinal_unbounded}.
\end{proof}

\begin{proposition}\label{positive_ordinal_with_unbounded_subset_is_limit}
    Let $\alpha$ be an ordinal.
    Suppose $\alpha$ is inhabited.
    Suppose $\ordinalUnbounded{C}{\alpha}$.
    Then $\alpha$ is a limit ordinal.
\end{proposition}
\begin{proof}
    $\alpha\neq\emptyset$ by \auto{nonempty_inhabited}.
    Hence $\emptyset\in\alpha$
        by \obv{ordinal_empty_or_emptyset_elem}.
    Show that for every $\beta\in\alpha$ we have $\suc{\beta}\in\alpha$.
    \begin{subproof}
        Fix $\beta\in\alpha$.
        Take $\gamma\in C$ such that $\beta\in\gamma$
            by \obv{ordinal_unbounded}.
        $C\subseteq\alpha$ by \obv{ordinal_unbounded}.
        Hence $\gamma\in\alpha$ by \obv{subseteq}.
        Hence $\gamma$ is an ordinal by \obv{elem_of_ordinal_is_ordinal}.
        Hence $\suc{\beta}\subseteq\gamma$
            by \obv{ordinal_suc_subseteq}.
        $\beta$ is an ordinal by \obv{elem_of_ordinal_is_ordinal}.
        Hence $\suc{\beta}$ is an ordinal by \obv{suc_ordinal}.
        Hence $\suc{\beta}\in\gamma$ or $\suc{\beta}=\gamma$
            by \obv{ordinal_subseteq_implies_elem_or_eq}.
        Follows by \obv{ordinal_transitivity}.
    \end{subproof}
    Hence $\emptyset\precedes\alpha$.
    Follows by \obv{limit_ordinal}.
\end{proof}

\begin{proposition}\label{ordinal_tail_is_club}
    Let $\alpha$ be a limit ordinal.
    Let $\beta\in\alpha$.
    Then $\clubIn{\ordinalTail{\beta}{\alpha}}{\alpha}$.
\end{proposition}
\begin{proof}
    Let $T=\ordinalTail{\beta}{\alpha}$.
    $T=\alpha\setminus\beta$ by \obv{ordinal_tail}.
    Hence $T\subseteq\alpha$ by \auto{setminus_subseteq}.
    Show that $\ordinalUnbounded{T}{\alpha}$.
    \begin{subproof}
        Show that for every $\delta\in\alpha$ there exists $\gamma\in T$
            such that $\delta\in\gamma$.
        \begin{subproof}
            Fix $\delta\in\alpha$.
            Hence $\delta$ and $\beta$ are ordinals
                by \obv{elem_of_ordinal_is_ordinal}.
            Let $m=\delta\union\beta$.
            $\delta\subseteq\beta$ or $\beta\subseteq\delta$
                by \obv{ordinal_subseteq_or_subseteq}.
            Hence $m=\delta$ or $m=\beta$
                by \obv{union_absorb_subseteq_left,union_absorb_subseteq_right}.
            Hence $m\in\alpha$.
            Let $\gamma=\suc{m}$.
            $\gamma\in\alpha$ by \auto{suc_elem_limitordinal}.
            $\delta\subseteq m$ and $\beta\subseteq m$
                by \obv{union_upper_left,union_upper_right}.
            Hence $\delta\in m$ or $\delta=m$
                by \obv{ordinal_subseteq_implies_elem_or_eq}.
            Hence $\delta\in\gamma$ by \auto{suc_iff}.
            Show that $\gamma\notin\beta$.
            \begin{subproof}
                Suppose $\gamma\in\beta$.
                $\beta\subseteq m$.
                Hence $\gamma\in m$ by \obv{subseteq}.
                $m\in\gamma$ by \obv{suc_intro_self}.
                Contradiction by \obv{in_asymmetric}.
            \end{subproof}
            Hence $\gamma\in\alpha\setminus\beta$
                by \obv{setminus_intro}.
            Hence $\gamma\in T$.
        \end{subproof}
        Follows by \obv{ordinal_unbounded}.
    \end{subproof}
    Show that $\ordinalClosed{T}{\alpha}$.
    \begin{subproof}
        Show that for every limit ordinal $\gamma\in\alpha$ such that
            $\ordinalUnbounded{T\inter\gamma}{\gamma}$
            we have $\gamma\in T$.
        \begin{subproof}
            Fix $\gamma$.
            Assume $\gamma\in\alpha$.
            Assume $\gamma$ is a limit ordinal.
            Assume $\ordinalUnbounded{T\inter\gamma}{\gamma}$.
            Show that $\gamma\notin\beta$.
            \begin{subproof}
                Suppose $\gamma\in\beta$.
                $\emptyset\in\gamma$ by \obv{limit_ordinal}.
                Take $\eta\in T\inter\gamma$ such that $\emptyset\in\eta$
                    by \obv{ordinal_unbounded}.
                Hence $\eta\in T$ and $\eta\in\gamma$
                    by \obv{inter_elim_left,inter_elim_right}.
                $\beta$ is an ordinal
                    by \obv{elem_of_ordinal_is_ordinal,limit_ordinal}.
                Hence $\eta\in\beta$ by \obv{ordinal_transitivity}.
                $\eta\notin\beta$ by \auto{ordinal_tail,setminus_elim_right}.
                Contradiction.
            \end{subproof}
            Hence $\gamma\in\alpha\setminus\beta$
                by \obv{setminus_intro}.
            Follows by \obv{ordinal_tail}.
        \end{subproof}
        Follows by \obv{ordinal_closed}.
    \end{subproof}
    Follows by \obv{club_in}.
\end{proof}

\begin{proposition}\label{ordinal_closed_intersection}
    Let $\alpha$ be an ordinal.
    Suppose $\ordinalClosed{C}{\alpha}$ and
        $\ordinalClosed{D}{\alpha}$.
    Then $\ordinalClosed{C\inter D}{\alpha}$.
\end{proposition}
\begin{proof}
    $C\subseteq\alpha$ and $D\subseteq\alpha$
        by \obv{ordinal_closed}.
    Hence $C\inter D\subseteq\alpha$ by \obv{subseteq,inter}.
    Show that for every limit ordinal $\gamma\in\alpha$ such that
        $\ordinalUnbounded{(C\inter D)\inter\gamma}{\gamma}$
        we have $\gamma\in C\inter D$.
    \begin{subproof}
        Fix $\gamma$.
        Assume $\gamma\in\alpha$.
        Assume $\gamma$ is a limit ordinal.
        Assume $\ordinalUnbounded{(C\inter D)\inter\gamma}{\gamma}$.
        $(C\inter D)\inter\gamma\subseteq C\inter\gamma$ and
            $(C\inter D)\inter\gamma\subseteq D\inter\gamma$
            by \auto{subseteq,inter}.
        $C\inter\gamma\subseteq\gamma$ and $D\inter\gamma\subseteq\gamma$
            by \obv{subseteq,inter}.
        Hence $\ordinalUnbounded{C\inter\gamma}{\gamma}$ and
            $\ordinalUnbounded{D\inter\gamma}{\gamma}$
            by \obv{ordinal_unbounded_superset}.
        Hence $\gamma\in C$ and $\gamma\in D$
            by \obv{ordinal_closed}.
        Follows by \obv{inter_intro}.
    \end{subproof}
    Follows by \obv{ordinal_closed}.
\end{proof}

\begin{proposition}\label{stationary_is_unbounded}
    Suppose $S$ is stationary in $\alpha$.
    Then $\ordinalUnbounded{S}{\alpha}$.
\end{proposition}
\begin{proof}
    $\alpha$ is a limit ordinal and $S\subseteq\alpha$
        by \obv{stationary_in_elim}.
    Show that for every $\beta\in\alpha$ there exists $\gamma\in S$
        such that $\beta\in\gamma$.
    \begin{subproof}
        Fix $\beta\in\alpha$.
        $\suc{\beta}\in\alpha$ by \obv{suc_elem_limitordinal}.
        $\clubIn{\ordinalTail{\suc{\beta}}{\alpha}}{\alpha}$
            by \obv{ordinal_tail_is_club}.
        Hence $S\meets\ordinalTail{\suc{\beta}}{\alpha}$
            by \obv{stationary_in}.
        Take $\gamma$ such that $\gamma\in S$ and
            $\gamma\in\ordinalTail{\suc{\beta}}{\alpha}$
            by \auto{disjoint}.
        Hence $\gamma\in\alpha$ by \obv{subseteq}.
        $\gamma\notin\suc{\beta}$
            by \auto{ordinal_tail,setminus_elim_right}.
        Hence $\beta$ and $\gamma$ are ordinals
            by \obv{elem_of_ordinal_is_ordinal}.
        $\gamma\notin\beta$ and $\gamma\neq\beta$
            by \auto{suc_iff}.
        Hence $\beta\in\gamma$ by \obv{ordinal_elem_connex}.
    \end{subproof}
    Follows by \obv{ordinal_unbounded}.
\end{proof}

\begin{proposition}\label{stationary_superset}
    Suppose $S$ is stationary in $\alpha$.
    Suppose $S\subseteq T\subseteq\alpha$.
    Then $T$ is stationary in $\alpha$.
\end{proposition}
\begin{proof}
    Show that for every $C$ such that $\clubIn{C}{\alpha}$
        we have $T\meets C$.
    \begin{subproof}
        Fix $C$.
        Assume $\clubIn{C}{\alpha}$.
        Hence $S\meets C$ by \obv{stationary_in}.
        Take $x$ such that $x\in S$ and $x\in C$
            by \auto{disjoint}.
        Hence $x\in T$ by \obv{subseteq}.
        Follows by \auto{disjoint}.
    \end{subproof}
    Follows by \obv{stationary_in}.
\end{proof}

\begin{proposition}\label{bounded_is_nonstationary}
    Let $\alpha$ be a limit ordinal.
    Suppose $\ordinalBounded{B}{\alpha}$.
    Then $\nonstationaryIn{B}{\alpha}$.
\end{proposition}
\begin{proof}
    $B\subseteq\alpha$ by \obv{ordinal_bounded}.
    Take $\beta\in\alpha$ such that $B\subseteq\beta$
        by \obv{ordinal_bounded}.
    Let $T=\ordinalTail{\beta}{\alpha}$.
    $\clubIn{T}{\alpha}$ by \auto{ordinal_tail_is_club}.
    Show that $B\notmeets T$.
    \begin{subproof}
        Show that there exists no $x$ such that $x\in B,T$.
        \begin{subproof}
            Suppose there exists $x$ such that $x\in B,T$.
            Take $x$ such that $x\in B,T$.
            Hence $x\in\beta$ by \obv{subseteq}.
            $x\notin\beta$ by \auto{ordinal_tail,setminus_elim_right}.
            Contradiction.
        \end{subproof}
        Follows by \obv{disjoint}.
    \end{subproof}
    Follows by \obv{nonstationary_in}.
\end{proof}

\begin{proposition}\label{club_upper_inhabited}
    Suppose $\ordinalUnbounded{C}{\alpha}$.
    Let $\beta\in\alpha$.
    Then $\clubUpper{C}{\beta}{\alpha}$ is inhabited.
\end{proposition}
\begin{proof}
    Take $\gamma\in C$ such that $\beta\in\gamma$
        by \obv{ordinal_unbounded}.
    Hence $\gamma\in\clubUpper{C}{\beta}{\alpha}$
        by \obv{club_upper}.
\end{proof}

\begin{proposition}\label{club_upper_are_ordinals}
    Let $\alpha$ be an ordinal.
    Suppose $C\subseteq\alpha$.
    For every $\gamma\in\clubUpper{C}{\beta}{\alpha}$ we have
        $\gamma$ is an ordinal.
\end{proposition}
\begin{proof}
    Fix $\gamma\in\clubUpper{C}{\beta}{\alpha}$.
    Hence $\gamma\in C$ by \obv{club_upper}.
    Hence $\gamma\in\alpha$ by \obv{subseteq}.
    Follows by \obv{elem_of_ordinal_is_ordinal}.
\end{proof}

\begin{proposition}\label{least_member_above_properties}
    Let $\alpha$ be an ordinal.
    Suppose $\ordinalUnbounded{C}{\alpha}$.
    Let $\beta\in\alpha$.
    Then $\leastMemberAbove{C}{\beta}{\alpha}\in C$ and
        $\beta\in\leastMemberAbove{C}{\beta}{\alpha}$ and
        $\leastMemberAbove{C}{\beta}{\alpha}\in\alpha$ and
        for every $\gamma\in C$ such that $\beta\in\gamma$ we have
            $\leastMemberAbove{C}{\beta}{\alpha}\subseteq\gamma$.
\end{proposition}
\begin{proof}
    Let $U=\clubUpper{C}{\beta}{\alpha}$.
    $C\subseteq\alpha$ by \obv{ordinal_unbounded}.
    $U$ is inhabited by \auto{club_upper_inhabited}.
    For every $\gamma\in U$ we have $\gamma$ is an ordinal
        by \auto{club_upper_are_ordinals}.
    Hence $\inters{U}\in U$ by \obv{inters_of_ordinals_elem}.
    Hence $\inters{U}\in C$ and $\beta\in\inters{U}$
        by \obv{club_upper}.
    Hence $\inters{U}\in\alpha$ by \obv{subseteq}.
    For every $\gamma\in C$ such that $\beta\in\gamma$ we have
        $\inters{U}\subseteq\gamma$
        by \obv{club_upper,inters_of_ordinals_subseteq}.
    Follows by \auto{least_member_above}.
\end{proof}

\begin{proposition}\label{least_member_above_candidates_inhabited}
    Suppose $\ordinalUnbounded{C}{\alpha}$.
    Let $\beta\in\alpha$.
    Then $\clubUpper{C}{\beta}{\alpha}$ is inhabited.
\end{proposition}
\begin{proof}
    Follows by \obv{club_upper_inhabited}.
\end{proof}

\begin{proposition}\label{least_member_above_mem}
    Let $\alpha$ be an ordinal.
    Suppose $\ordinalUnbounded{C}{\alpha}$.
    Let $\beta\in\alpha$.
    Then $\leastMemberAbove{C}{\beta}{\alpha}\in C$.
\end{proposition}
\begin{proof}
    Follows by \obv{least_member_above_properties}.
\end{proof}

\begin{proposition}\label{least_member_above_strict}
    Let $\alpha$ be an ordinal.
    Suppose $\ordinalUnbounded{C}{\alpha}$.
    Let $\beta\in\alpha$.
    Then $\beta\in\leastMemberAbove{C}{\beta}{\alpha}$.
\end{proposition}
\begin{proof}
    Follows by \obv{least_member_above_properties}.
\end{proof}

\begin{proposition}\label{least_member_above_below_ambient}
    Let $\alpha$ be an ordinal.
    Suppose $\ordinalUnbounded{C}{\alpha}$.
    Let $\beta\in\alpha$.
    Then $\leastMemberAbove{C}{\beta}{\alpha}\in\alpha$.
\end{proposition}
\begin{proof}
    Follows by \obv{least_member_above_properties}.
\end{proof}

\begin{proposition}\label{least_member_above_least}
    Let $\alpha$ be an ordinal.
    Suppose $\ordinalUnbounded{C}{\alpha}$.
    Let $\beta\in\alpha$.
    For every $\gamma\in C$ such that $\beta\in\gamma$ we have
        $\leastMemberAbove{C}{\beta}{\alpha}\subseteq\gamma$.
\end{proposition}
\begin{proof}
    Follows by \obv{least_member_above_properties}.
\end{proof}

\begin{proposition}\label{least_member_above_unique}
    Let $\alpha$ be an ordinal.
    Suppose $\ordinalUnbounded{C}{\alpha}$.
    Let $\beta\in\alpha$.
    Let $\eta\in C$.
    Suppose $\beta\in\eta$.
    Suppose for every $\gamma\in C$ such that $\beta\in\gamma$ we have
        $\eta\subseteq\gamma$.
    Then $\eta=\leastMemberAbove{C}{\beta}{\alpha}$.
\end{proposition}
\begin{proof}
    $\leastMemberAbove{C}{\beta}{\alpha}\subseteq\eta$
        by \obv{least_member_above_properties}.
    $\leastMemberAbove{C}{\beta}{\alpha}\in C$ and
        $\beta\in\leastMemberAbove{C}{\beta}{\alpha}$
        by \obv{least_member_above_properties}.
    Hence $\eta\subseteq\leastMemberAbove{C}{\beta}{\alpha}$.
    Follows by \obv{subseteq_antisymmetric}.
\end{proof}

\begin{theorem}\label{closed_club_gap_straddles}
    Suppose $\clubIn{C}{\beta}$.
    Let $\alpha\in\beta$.
    Suppose $\alpha$ is a limit ordinal and $\alpha\notin C$.
    Let $\xi=\unions{C\inter\alpha}$.
    Then $\xi\in\alpha$ and
        $\alpha\in\leastMemberAbove{C}{\xi}{\beta}$ and
        $\leastMemberAbove{C}{\xi}{\beta}\in\beta$.
\end{theorem}
\begin{proof}
    $\beta$ is a limit ordinal and $C\subseteq\beta$ and
        $\ordinalUnbounded{C}{\beta}$ and $\ordinalClosed{C}{\beta}$
        by \obv{club_in_elim}.
    Hence $\beta$ is an ordinal by \obv{limit_ordinal}.
    $C\inter\alpha\subseteq\alpha$ by \obv{subseteq,inter}.
    Every element of $C\inter\alpha$ is an ordinal
        by \obv{inter_elim_right,elem_of_ordinal_is_ordinal,limit_ordinal}.
    Hence $\xi$ is an ordinal by \obv{union_of_ordinals_is_ordinal}.
    Show that every element of $\xi$ is an element of $\alpha$.
    \begin{subproof}
        Fix $x\in\xi$.
        Take $\gamma\in C\inter\alpha$ such that $x\in\gamma$
            by \obv{unions_iff}.
        Hence $\gamma\in\alpha$ by \obv{inter_elim_right}.
        Follows by \obv{ordinal_transitivity,limit_ordinal}.
    \end{subproof}
    Hence $\xi\subseteq\alpha$ by \obv{subseteq}.
    Show that $\xi\neq\alpha$.
    \begin{subproof}
        Suppose $\xi=\alpha$.
        Hence $\ordinalUnbounded{C\inter\alpha}{\alpha}$
            by \auto{ordinal_unbounded_iff_union}.
        Hence $\alpha\in C$ by \obv{ordinal_closed}.
        Contradiction.
    \end{subproof}
    Hence $\xi\in\alpha$
        by \obv{ordinal_subseteq_implies_elem_or_eq,limit_ordinal}.
    Hence $\xi\in\beta$ by \obv{ordinal_transitivity}.
    Let $\eta=\leastMemberAbove{C}{\xi}{\beta}$.
    $\eta\in C$ and $\xi\in\eta$ and $\eta\in\beta$
        by \auto{least_member_above_properties}.
    Hence $\eta$ is an ordinal
        by \obv{elem_of_ordinal_is_ordinal}.
    Show that $\alpha\in\eta$.
    \begin{subproof}
        Suppose not.
        Hence $\eta\in\alpha$ or $\alpha=\eta$
            by \obv{ordinal_elem_connex,limit_ordinal}.
        $\alpha\neq\eta$.
        Hence $\eta\in\alpha$.
        Hence $\eta\in C\inter\alpha$ by \obv{inter_intro}.
        Hence $\eta\subseteq\xi$ by \auto{subseteq,unions_iff}.
        Hence $\xi\in\xi$ by \obv{subseteq}.
        Contradiction by \obv{in_irrefl}.
    \end{subproof}
\end{proof}

\begin{proposition}\label{club_intermediate_properties}
    Suppose $\ordinalUnbounded{C}{\alpha}$.
    Let $\beta\in\alpha$.
    Then $\clubIntermediate{C}{\beta}{\alpha}\in C$ and
        $\beta\in\clubIntermediate{C}{\beta}{\alpha}$.
\end{proposition}
\begin{proof}
    $\clubUpper{C}{\beta}{\alpha}$ is inhabited
        by \obv{club_upper_inhabited}.
    Hence $\choice{\clubUpper{C}{\beta}{\alpha}}
        \in\clubUpper{C}{\beta}{\alpha}$ by \obv{choice_in}.
    Follows by \obv{club_intermediate,club_upper}.
\end{proof}

\begin{theorem}\label{club_intersection_advance_type}
    Suppose $\ordinalUnbounded{C}{\alpha}$ and
        $\ordinalUnbounded{D}{\alpha}$.
    Then $\clubIntersectionAdvance{C}{D}{\alpha}
        \in\funs{\alpha}{\alpha}$.
\end{theorem}
\begin{proof}
    Let $f=\clubIntersectionAdvance{C}{D}{\alpha}$.
    $C\subseteq\alpha$ and $D\subseteq\alpha$
        by \obv{ordinal_unbounded}.
    $f$ is a relation by \auto{relation,club_intersection_advance}.
    $f$ is right-unique
        by \auto{rightunique,club_intersection_advance,pair_eq_iff}.
    $\dom{f}=\alpha$
        by \auto{setext,dom_iff,club_intersection_advance,pair_eq_iff}.
    Show that for every $\beta\in\dom{f}$ we have $f(\beta)\in\alpha$.
    \begin{subproof}
        Fix $\beta\in\dom{f}$.
        Hence $\beta\in\alpha$.
        $\clubIntermediate{C}{\beta}{\alpha}\in C$
            by \obv{club_intermediate_properties}.
        Hence $\clubIntermediate{C}{\beta}{\alpha}\in\alpha$
            by \obv{subseteq}.
        $\clubUpper{D}{\clubIntermediate{C}{\beta}{\alpha}}{\alpha}$
            is inhabited by \obv{club_upper_inhabited}.
        Hence $\choice{\clubUpper{D}{
            \clubIntermediate{C}{\beta}{\alpha}}{\alpha}}
            \in\clubUpper{D}{
                \clubIntermediate{C}{\beta}{\alpha}}{\alpha}$
            by \obv{choice_in}.
        Hence $\choice{\clubUpper{D}{
            \clubIntermediate{C}{\beta}{\alpha}}{\alpha}}\in D$
            by \obv{club_upper}.
        Hence $\choice{\clubUpper{D}{
            \clubIntermediate{C}{\beta}{\alpha}}{\alpha}}\in\alpha$
            by \obv{subseteq}.
        $(\beta,\choice{\clubUpper{D}{
            \clubIntermediate{C}{\beta}{\alpha}}{\alpha}})\in f$
            by \obv{club_intersection_advance}.
        Hence $f(\beta)=\choice{\clubUpper{D}{
            \clubIntermediate{C}{\beta}{\alpha}}{\alpha}}$
            by \obv{function_apply_intro}.
    \end{subproof}
    Hence $f$ is a function to $\alpha$.
    Follows by \auto{funs_intro}.
\end{proof}

\begin{proposition}\label{club_intersection_advance_properties}
    Let $\alpha$ be an ordinal.
    Suppose $\ordinalUnbounded{C}{\alpha}$ and
        $\ordinalUnbounded{D}{\alpha}$.
    Let $\beta\in\alpha$.
    Let $f=\clubIntersectionAdvance{C}{D}{\alpha}$.
    Then $f(\beta)\in D$ and
        $\clubIntermediate{C}{\beta}{\alpha}\in f(\beta)$ and
        $\beta\in f(\beta)$.
\end{proposition}
\begin{proof}
    $C\subseteq\alpha$ by \obv{ordinal_unbounded}.
    $\clubIntermediate{C}{\beta}{\alpha}\in C$ and
        $\beta\in\clubIntermediate{C}{\beta}{\alpha}$
        by \obv{club_intermediate_properties}.
    Hence $\clubIntermediate{C}{\beta}{\alpha}\in\alpha$
        by \obv{subseteq}.
    $\clubUpper{D}{\clubIntermediate{C}{\beta}{\alpha}}{\alpha}$
        is inhabited by \obv{club_upper_inhabited}.
    Hence $\choice{\clubUpper{D}{
        \clubIntermediate{C}{\beta}{\alpha}}{\alpha}}
        \in\clubUpper{D}{
            \clubIntermediate{C}{\beta}{\alpha}}{\alpha}$
        by \obv{choice_in}.
    Hence $\choice{\clubUpper{D}{
        \clubIntermediate{C}{\beta}{\alpha}}{\alpha}}\in D$ and
        $\clubIntermediate{C}{\beta}{\alpha}
            \in\choice{\clubUpper{D}{
                \clubIntermediate{C}{\beta}{\alpha}}{\alpha}}$
        by \obv{club_upper}.
    $f$ is a function
        by \auto{club_intersection_advance_type,funs_is_function}.
    $(\beta,\choice{\clubUpper{D}{
        \clubIntermediate{C}{\beta}{\alpha}}{\alpha}})\in f$
        by \obv{club_intersection_advance}.
    Hence $f(\beta)=\choice{\clubUpper{D}{
        \clubIntermediate{C}{\beta}{\alpha}}{\alpha}}$
        by \obv{function_apply_intro}.
    Hence $f(\beta)\in D$.
    $D\subseteq\alpha$ by \obv{ordinal_unbounded}.
    Hence $f(\beta)\in\alpha$ by \obv{subseteq}.
    $f(\beta)$ is an ordinal
        by \obv{elem_of_ordinal_is_ordinal}.
    Hence $\beta\in f(\beta)$ by \obv{ordinal_transitivity}.
\end{proof}

\begin{proposition}\label{club_intersection_iteration_properties}
    Let $\alpha$ be an ordinal.
    Suppose $\ordinalUnbounded{C}{\alpha}$ and
        $\ordinalUnbounded{D}{\alpha}$.
    Let $\beta\in\alpha$.
    Let $f=\clubIntersectionAdvance{C}{D}{\alpha}$.
    Let $a=f(\beta)$.
    Let $s=\iterationSet{f}{a}$.
    Then $s\in\funs{\naturals}{\alpha}$ and
        for every natural number $n$ we have
            $s(n)\in D$ and $s(n)\in s(\suc{n})$ and
            $\clubIntermediate{C}{s(n)}{\alpha}\in s(\suc{n})$.
\end{proposition}
\begin{proof}
    $f\in\funs{\alpha}{\alpha}$
        by \auto{club_intersection_advance_type}.
    $a\in\alpha$ by \auto{funs_type_apply}.
    $s\in\funs{\naturals}{\alpha}$
        by \auto{iteration_set_type}.
    Hence $s$ is a function and $\dom{s}=\naturals$
        by \obv{funs_is_function,funs}.
    Let $N=\{n\in\naturals\mid
        s(n)\in D\land s(n)\in s(\suc{n})\land
        \clubIntermediate{C}{s(n)}{\alpha}\in s(\suc{n})\}$.
    Show that $\emptyset\in N$.
    \begin{subproof}
        $s(\emptyset)=a$ by \auto{iteration_set_zero}.
        $a\in D$ by \auto{club_intersection_advance_properties}.
        $\emptyset$ is a natural number by \obv{omega_empty}.
        $s(\suc{\emptyset})=f(s(\emptyset))$
            by \auto{iteration_set_successor}.
        $s(\emptyset)\in\alpha$ by \obv{funs_type_apply}.
        $f(s(\emptyset))\in D$ and
            $\clubIntermediate{C}{s(\emptyset)}{\alpha}
                \in f(s(\emptyset))$ and
            $s(\emptyset)\in f(s(\emptyset))$
            by \obv{club_intersection_advance_properties}.
    \end{subproof}
    Show that for every $n\in N$ we have $\suc{n}\in N$.
    \begin{subproof}
        Fix $n\in N$.
        Hence $\suc{n}$ is a natural number by \obv{omega_successor}.
        $s(n)\in\alpha$ by \obv{funs_type_apply}.
        $s(\suc{n})=f(s(n))$
            by \auto{iteration_set_successor}.
        $f(s(n))\in D$ by
            \obv{club_intersection_advance_properties}.
        Hence $s(\suc{n})\in D$.
        $s(\suc{n})\in\alpha$ by \obv{funs_type_apply}.
        $s(\suc{\suc{n}})=f(s(\suc{n}))$
            by \auto{iteration_set_successor}.
        $f(s(\suc{n}))\in D$ and
            $\clubIntermediate{C}{s(\suc{n})}{\alpha}
                \in f(s(\suc{n}))$ and
            $s(\suc{n})\in f(s(\suc{n}))$
            by \obv{club_intersection_advance_properties}.
        Hence $s(\suc{n})\in s(\suc{\suc{n}})$.
        $\clubIntermediate{C}{s(\suc{n})}{\alpha}
            \in s(\suc{\suc{n}})$.
    \end{subproof}
    Hence $N$ is an inductive set by \obv{inductive_set}.
    Hence every natural number is an element of $N$
        by \obv{omega_minimal}.
    Show that for every natural number $n$ we have
        $s(n)\in D$ and $s(n)\in s(\suc{n})$ and
        $\clubIntermediate{C}{s(n)}{\alpha}\in s(\suc{n})$.
    \begin{subproof}
        Fix $n\in\naturals$.
        $n\in N$.
    \end{subproof}
\end{proof}

\begin{theorem}\label{two_clubs_have_common_limit_point_above}
    Let $\alpha$ be a limit ordinal.
    Suppose $\naturals\in\cofinality{\alpha}$.
    Suppose $\clubIn{C}{\alpha}$ and $\clubIn{D}{\alpha}$.
    Let $\beta\in\alpha$.
    Then there exists $x\in\alpha$ such that $\beta\in x$ and
        $x$ is a limit ordinal and
        $\ordinalUnbounded{C\inter x}{x}$ and
        $\ordinalUnbounded{D\inter x}{x}$.
\end{theorem}
\begin{proof}
    $\alpha$ is an ordinal by \obv{limit_ordinal}.
    $\ordinalUnbounded{C}{\alpha}$ and
        $\ordinalClosed{C}{\alpha}$ and
        $\ordinalUnbounded{D}{\alpha}$ and
        $\ordinalClosed{D}{\alpha}$
        by \obv{club_in_elim}.
    Let $f=\clubIntersectionAdvance{C}{D}{\alpha}$.
    $f\in\funs{\alpha}{\alpha}$
        by \auto{club_intersection_advance_type}.
    Let $a=f(\beta)$.
    $a\in\alpha$ by \auto{funs_type_apply}.
    Let $s=\iterationSet{f}{a}$.
    $s\in\funs{\naturals}{\alpha}$ and
        for every natural number $n$ we have
            $s(n)\in D$ and $s(n)\in s(\suc{n})$ and
            $\clubIntermediate{C}{s(n)}{\alpha}\in s(\suc{n})$
        by \obv{club_intersection_iteration_properties}.
    Hence $s$ is a function and $\dom{s}=\naturals$
        by \obv{funs_is_function,funs}.
    Let $x=\unions{\ran{s}}$.
    $x\in\alpha$ by \auto{small_domain_range_union_below}.
    $\ran{s}\subseteq\alpha$ by \obv{funs_ran}.
    Hence every element of $\ran{s}$ is an ordinal
        by \obv{subseteq,elem_of_ordinal_is_ordinal}.
    Hence $x$ is an ordinal by \obv{union_of_ordinals_is_ordinal}.
    $\emptyset$ is a natural number by \obv{omega_empty}.
    Hence $s(\emptyset)\in s(\suc{\emptyset})$
        by \auto{club_intersection_iteration_properties}.
    $\suc{\emptyset}$ is a natural number by \obv{omega_successor}.
    Hence $s(\suc{\emptyset})\in\ran{s}$ by \auto{fun_ran_iff}.
    Hence $s(\emptyset)\in x$ by \auto{unions_iff}.
    Hence $x\neq\emptyset$ by \auto{emptyset}.
    Hence $\emptyset\in x$
        by \obv{ordinal_empty_or_emptyset_elem}.
    Show that for every $\rho\in x$ we have $\suc{\rho}\in x$.
    \begin{subproof}
        Fix $\rho\in x$.
        Take $y\in\ran{s}$ such that $\rho\in y$
            by \obv{unions_iff}.
        Take $n\in\dom{s}$ such that $s(n)=y$
            by \obv{fun_ran_iff}.
        Hence $n$ is a natural number.
        Hence $s(n)\in s(\suc{n})$
            by \auto{club_intersection_iteration_properties}.
        $s(n)$ is an ordinal
            by \obv{funs_type_apply,elem_of_ordinal_is_ordinal}.
        Hence $\rho$ is an ordinal by \obv{elem_of_ordinal_is_ordinal}.
        Hence $\suc{\rho}\subseteq s(n)$
            by \obv{ordinal_suc_subseteq}.
        $\suc{\rho}$ is an ordinal by \obv{suc_ordinal}.
        Hence $\suc{\rho}\in s(n)$ or $\suc{\rho}=s(n)$
            by \obv{ordinal_subseteq_implies_elem_or_eq}.
        $s(\suc{n})$ is an ordinal
            by \obv{funs_type_apply,omega_successor,elem_of_ordinal_is_ordinal}.
        Hence $\suc{\rho}\in s(\suc{n})$
            by \obv{ordinal_transitivity}.
        $\suc{n}$ is a natural number by \obv{omega_successor}.
        Hence $s(\suc{n})\in\ran{s}$ by \auto{fun_ran_iff}.
        Follows by \auto{unions_iff}.
    \end{subproof}
    Hence $\emptyset\precedes x$.
    Hence $x$ is a limit ordinal by \obv{limit_ordinal}.
    Show that $\ordinalUnbounded{C\inter x}{x}$.
    \begin{subproof}
        $C\inter x\subseteq x$ by \obv{subseteq,inter}.
        Show that for every $\rho\in x$ there exists $m\in C\inter x$
            such that $\rho\in m$.
        \begin{subproof}
            Fix $\rho\in x$.
            Take $y\in\ran{s}$ such that $\rho\in y$
                by \obv{unions_iff}.
            Take $n\in\dom{s}$ such that $s(n)=y$
                by \obv{fun_ran_iff}.
            Hence $n$ is a natural number.
            Let $m=\clubIntermediate{C}{s(n)}{\alpha}$.
            $s(n)\in\alpha$ by \obv{funs_type_apply}.
            Hence $m\in C$ and $s(n)\in m$
                by \auto{club_intermediate_properties}.
            $m\in s(\suc{n})$
                by \auto{club_intersection_iteration_properties}.
            $\suc{n}$ is a natural number by \obv{omega_successor}.
            Hence $s(\suc{n})\in\ran{s}$ by \auto{fun_ran_iff}.
            Hence $m\in x$ by \auto{unions_iff}.
            Hence $m\in C\inter x$ by \obv{inter_intro}.
            $C\subseteq\alpha$ by \obv{ordinal_unbounded}.
            Hence $m\in\alpha$ by \obv{subseteq}.
            Hence $m$ is an ordinal by \obv{elem_of_ordinal_is_ordinal}.
            $\rho\in s(n)$.
            Hence $\rho\in m$ by \obv{ordinal_transitivity}.
        \end{subproof}
        Follows by \obv{ordinal_unbounded}.
    \end{subproof}
    Show that $\ordinalUnbounded{D\inter x}{x}$.
    \begin{subproof}
        $D\inter x\subseteq x$ by \obv{subseteq,inter}.
        Show that for every $\rho\in x$ there exists $z\in D\inter x$
            such that $\rho\in z$.
        \begin{subproof}
            Fix $\rho\in x$.
            Take $y\in\ran{s}$ such that $\rho\in y$
                by \obv{unions_iff}.
            Take $n\in\dom{s}$ such that $s(n)=y$
                by \obv{fun_ran_iff}.
            Hence $n$ is a natural number.
            $s(\suc{n})\in D$
                by \auto{club_intersection_iteration_properties,omega_successor}.
            $s(n)\in s(\suc{n})$
                by \auto{club_intersection_iteration_properties}.
            Let $z=s(\suc{n})$.
            Hence $z\in D$.
            $s(n)\in z$.
            $\suc{n}$ is a natural number by \obv{omega_successor}.
            Hence $z\in s(\suc{\suc{n}})$
                by \auto{club_intersection_iteration_properties,omega_successor}.
            $\suc{\suc{n}}$ is a natural number
                by \obv{omega_successor}.
            Hence $s(\suc{\suc{n}})\in\ran{s}$
                by \auto{fun_ran_iff}.
            Hence $z\in x$ by \auto{unions_iff}.
            Hence $z\in D\inter x$ by \obv{inter_intro}.
            $D\subseteq\alpha$ by \obv{ordinal_unbounded}.
            Hence $z\in\alpha$ by \obv{subseteq}.
            Hence $z$ is an ordinal by \obv{elem_of_ordinal_is_ordinal}.
            $\rho\in s(n)$.
            Hence $\rho\in z$ by \auto{ordinal_transitivity}.
            Follows by \auto.
        \end{subproof}
        Follows by \obv{ordinal_unbounded}.
    \end{subproof}
    Hence $x\in C$ and $x\in D$ by \obv{ordinal_closed}.
    Hence $x\in C\inter D$ by \obv{inter_intro}.
    $\beta\in a$ by \auto{club_intersection_advance_properties}.
    $s(\emptyset)=a$ by \auto{iteration_set_zero}.
    Hence $\beta\in s(\emptyset)$.
    Hence $\beta\in x$ by \auto{ordinal_transitivity}.
    Follows by \auto.
\end{proof}

\begin{theorem}\label{two_clubs_have_common_point_above}
    Let $\alpha$ be a limit ordinal.
    Suppose $\naturals\in\cofinality{\alpha}$.
    Suppose $\clubIn{C}{\alpha}$ and $\clubIn{D}{\alpha}$.
    Let $\beta\in\alpha$.
    Then there exists $x\in C\inter D$ such that
        $\beta\in x$ and $x$ is a limit ordinal.
\end{theorem}
\begin{proof}
    Take $x\in\alpha$ such that $\beta\in x$ and
        $x$ is a limit ordinal and
        $\ordinalUnbounded{C\inter x}{x}$ and
        $\ordinalUnbounded{D\inter x}{x}$
        by \obv{two_clubs_have_common_limit_point_above}.
    $\ordinalClosed{C}{\alpha}$ and $\ordinalClosed{D}{\alpha}$
        by \obv{club_in_elim}.
    Hence $x\in C$ and $x\in D$ by \obv{ordinal_closed}.
    Hence $x\in C\inter D$ by \obv{inter_intro}.
    Follows by \auto.
\end{proof}

\begin{theorem}\label{two_club_intersection}
    Let $\alpha$ be a limit ordinal.
    Suppose $\naturals\in\cofinality{\alpha}$.
    Suppose $\clubIn{C}{\alpha}$ and $\clubIn{D}{\alpha}$.
    Then $\clubIn{C\inter D}{\alpha}$.
\end{theorem}
\begin{proof}
    $C\subseteq\alpha$ and $D\subseteq\alpha$
        by \auto{club_in_elim}.
    Hence $C\inter D\subseteq\alpha$ by \obv{subseteq,inter}.
    Show that $\ordinalUnbounded{C\inter D}{\alpha}$.
    \begin{subproof}
        For every $\beta\in\alpha$ there exists $x\in C\inter D$ such that
            $\beta\in x$ by \auto{two_clubs_have_common_point_above}.
        Follows by \obv{ordinal_unbounded}.
    \end{subproof}
    $\ordinalClosed{C}{\alpha}$ and $\ordinalClosed{D}{\alpha}$
        by \obv{club_in_elim}.
    Hence $\ordinalClosed{C\inter D}{\alpha}$
        by \obv{ordinal_closed_intersection}.
    Follows by \obv{club_in}.
\end{proof}

\begin{proposition}\label{club_is_stationary}
    Let $\alpha$ be a limit ordinal.
    Suppose $\naturals\in\cofinality{\alpha}$.
    Suppose $\clubIn{C}{\alpha}$.
    Then $C$ is stationary in $\alpha$.
\end{proposition}
\begin{proof}
    $C\subseteq\alpha$ by \auto{club_in_elim}.
    Show that for every $D$ such that $\clubIn{D}{\alpha}$
        we have $C\meets D$.
    \begin{subproof}
        Fix $D$.
        Assume $\clubIn{D}{\alpha}$.
        $\emptyset\in\alpha$ by \obv{limit_ordinal}.
        Take $x\in C\inter D$ such that $\emptyset\in x$
            by \auto{two_clubs_have_common_point_above}.
        Hence $x\in C$ and $x\in D$
            by \obv{inter_elim_left,inter_elim_right}.
        Follows by \auto{disjoint}.
    \end{subproof}
    Follows by \obv{stationary_in}.
\end{proof}

\begin{proposition}\label{nonstationary_iff_not_stationary}
    Let $\alpha$ be a limit ordinal.
    Suppose $S\subseteq\alpha$.
    Then $\nonstationaryIn{S}{\alpha}$ iff
        $S$ is not stationary in $\alpha$.
\end{proposition}
\begin{proof}
    Show that if $\nonstationaryIn{S}{\alpha}$, then
        $S$ is not stationary in $\alpha$.
    \begin{subproof}
        Assume $\nonstationaryIn{S}{\alpha}$.
        Take $C$ such that $\clubIn{C}{\alpha}$ and $S\notmeets C$
            by \obv{nonstationary_in}.
        Suppose $S$ is stationary in $\alpha$.
        Hence $S\meets C$ by \obv{stationary_in}.
        Contradiction.
    \end{subproof}
    Show that if $S$ is not stationary in $\alpha$, then
        $\nonstationaryIn{S}{\alpha}$.
    \begin{subproof}
        Assume $S$ is not stationary in $\alpha$.
        Take $C$ such that $\clubIn{C}{\alpha}$ and
            $\lnot(S\meets C)$ by \auto{stationary_in}.
        Hence $S\notmeets C$.
        Follows by \auto{nonstationary_in}.
    \end{subproof}
    Follows by \auto.
\end{proof}

\begin{proposition}\label{nonstationary_union}
    Let $\alpha$ be a limit ordinal.
    Suppose $\naturals\in\cofinality{\alpha}$.
    Suppose $\nonstationaryIn{S}{\alpha}$ and
        $\nonstationaryIn{T}{\alpha}$.
    Then $\nonstationaryIn{S\union T}{\alpha}$.
\end{proposition}
\begin{proof}
    $S\subseteq\alpha$ and $T\subseteq\alpha$
        by \obv{nonstationary_in}.
    Hence $S\union T\subseteq\alpha$ by \auto{subseteq_union_iff}.
    Take $C$ such that $\clubIn{C}{\alpha}$ and $S\notmeets C$
        by \obv{nonstationary_in}.
    Take $D$ such that $\clubIn{D}{\alpha}$ and $T\notmeets D$
        by \obv{nonstationary_in}.
    $\clubIn{C\inter D}{\alpha}$ by \obv{two_club_intersection}.
    Show that $(S\union T)\notmeets(C\inter D)$.
    \begin{subproof}
        Show that there exists no $x$ such that
            $x\in S\union T,C\inter D$.
        \begin{subproof}
            Suppose there exists $x$ such that
                $x\in S\union T,C\inter D$.
            Take $x$ such that $x\in S\union T,C\inter D$.
            Hence $x\in S$ or $x\in T$ by \auto{union_iff}.
            $x\in C$ and $x\in D$
                by \obv{inter_elim_left,inter_elim_right}.
            Contradiction by \auto{disjoint}.
        \end{subproof}
        Follows by \obv{disjoint}.
    \end{subproof}
    Follows by \auto{nonstationary_in}.
\end{proof}

\begin{proposition}\label{nonstationary_subset}
    Suppose $\nonstationaryIn{N}{\alpha}$ and $A\subseteq N$.
    Then $\nonstationaryIn{A}{\alpha}$.
\end{proposition}
\begin{proof}
    $\alpha$ is a limit ordinal and $N\subseteq\alpha$
        by \obv{nonstationary_in}.
    Hence $A\subseteq\alpha$ by \obv{subseteq_transitive}.
    Take $C$ such that $\clubIn{C}{\alpha}$ and $N\notmeets C$
        by \obv{nonstationary_in}.
    Show that $A\notmeets C$.
    \begin{subproof}
        Suppose not.
        Take $x$ such that $x\in A,C$ by \auto{disjoint}.
        Hence $x\in N$ by \obv{subseteq}.
        Contradiction by \auto{disjoint}.
    \end{subproof}
    Follows by \auto{nonstationary_in}.
\end{proof}

\begin{theorem}\label{stationary_remove_nonstationary}
    Let $\alpha$ be a limit ordinal.
    Suppose $\naturals\in\cofinality{\alpha}$.
    Suppose $S$ is stationary in $\alpha$ and
        $\nonstationaryIn{N}{\alpha}$.
    Then $S\setminus N$ is stationary in $\alpha$.
\end{theorem}
\begin{proof}
    $S\subseteq\alpha$ by \obv{stationary_in_elim}.
    Hence $S\setminus N\subseteq\alpha$
        by \obv{setminus_subseteq,subseteq_transitive}.
    Take $C$ such that $\clubIn{C}{\alpha}$ and $N\notmeets C$
        by \obv{nonstationary_in}.
    Show that for every $D$ such that $\clubIn{D}{\alpha}$
        we have $(S\setminus N)\meets D$.
    \begin{subproof}
        Fix $D$.
        Assume $\clubIn{D}{\alpha}$.
        $\clubIn{C\inter D}{\alpha}$ by \obv{two_club_intersection}.
        Hence $S\meets(C\inter D)$ by \obv{stationary_in}.
        Take $x$ such that $x\in S$ and $x\in C\inter D$
            by \auto{disjoint}.
        Hence $x\in C$ and $x\in D$
            by \obv{inter_elim_left,inter_elim_right}.
        Hence $x\notin N$ by \auto{disjoint}.
        Hence $x\in S\setminus N$ by \obv{setminus_intro}.
        Follows by \auto{disjoint}.
    \end{subproof}
    Follows by \obv{stationary_in}.
\end{proof}

\begin{proposition}\label{stationary_union_component}
    Let $\alpha$ be a limit ordinal.
    Suppose $\naturals\in\cofinality{\alpha}$.
    Suppose $S\subseteq\alpha$.
    Suppose $T\subseteq\alpha$.
    Suppose $S\union T$ is stationary in $\alpha$.
    Then $S$ is stationary in $\alpha$ or $T$ is stationary in $\alpha$.
\end{proposition}
\begin{proof}
    Suppose not.
    Hence $S$ is not stationary in $\alpha$ and
        $T$ is not stationary in $\alpha$.
    Hence $\nonstationaryIn{S}{\alpha}$ and
        $\nonstationaryIn{T}{\alpha}$
        by \obv{nonstationary_iff_not_stationary}.
    Hence $\nonstationaryIn{S\union T}{\alpha}$
        by \obv{nonstationary_union}.
    $S\union T\subseteq\alpha$ by \obv{subseteq_union_iff}.
    Hence $S\union T$ is not stationary in $\alpha$
        by \obv{nonstationary_iff_not_stationary}.
    Contradiction.
\end{proof}

\begin{theorem}\label{stationary_remove_bounded}
    Let $\alpha$ be a limit ordinal.
    Suppose $\naturals\in\cofinality{\alpha}$.
    Suppose $S$ is stationary in $\alpha$ and
        $\ordinalBounded{B}{\alpha}$.
    Then $S\setminus B$ is stationary in $\alpha$.
\end{theorem}
\begin{proof}
    $S\subseteq\alpha$ by \obv{stationary_in_elim}.
    Hence $S\setminus B\subseteq\alpha$
        by \obv{setminus_subseteq,subseteq_transitive}.
    Take $\beta\in\alpha$ such that $B\subseteq\beta$
        by \obv{ordinal_bounded}.
    Let $T=\ordinalTail{\beta}{\alpha}$.
    $\clubIn{T}{\alpha}$ by \auto{ordinal_tail_is_club}.
    Show that for every $C$ such that $\clubIn{C}{\alpha}$
        we have $(S\setminus B)\meets C$.
    \begin{subproof}
        Fix $C$.
        Assume $\clubIn{C}{\alpha}$.
        $\clubIn{C\inter T}{\alpha}$
            by \obv{two_club_intersection}.
        Hence $S\meets(C\inter T)$ by \obv{stationary_in}.
        Take $x$ such that $x\in S$ and $x\in C\inter T$
            by \auto{disjoint}.
        Hence $x\in C$ and $x\in T$
            by \obv{inter_elim_left,inter_elim_right}.
        $x\notin\beta$ by \auto{ordinal_tail,setminus_elim_right}.
        Hence $x\notin B$ by \auto{subseteq}.
        Hence $x\in S\setminus B$ by \obv{setminus_intro}.
        Follows by \auto{disjoint}.
    \end{subproof}
    Follows by \obv{stationary_in}.
\end{proof}

\begin{corollary}\label{stationary_strict_tail}
    Let $\alpha$ be a limit ordinal.
    Suppose $\naturals\in\cofinality{\alpha}$.
    Suppose $S$ is stationary in $\alpha$.
    Let $\beta\in\alpha$.
    Then $S\setminus\suc{\beta}$ is stationary in $\alpha$.
\end{corollary}
\begin{proof}
    $\suc{\beta}\in\alpha$ by \obv{suc_elem_limitordinal}.
    $\beta$ is an ordinal by \obv{elem_of_ordinal_is_ordinal}.
    Hence $\suc{\beta}$ is an ordinal by \obv{suc_ordinal}.
    Hence $\suc{\beta}\subseteq\alpha$
        by \obv{ordinal_elem_implies_subseteq}.
    $\suc{\beta}\subseteq\suc{\beta}$ by \obv{subseteq_refl}.
    Hence there exists $\gamma\in\alpha$ such that
        $\suc{\beta}\subseteq\gamma$ by \auto.
    Hence $\ordinalBounded{\suc{\beta}}{\alpha}$
        by \auto{ordinal_bounded}.
    Follows by \obv{stationary_remove_bounded}.
\end{proof}

\begin{theorem}\label{stationary_intersection_club}
    Let $\alpha$ be a limit ordinal.
    Suppose $\naturals\in\cofinality{\alpha}$.
    Suppose $S$ is stationary in $\alpha$ and $\clubIn{C}{\alpha}$.
    Then $S\inter C$ is stationary in $\alpha$.
\end{theorem}
\begin{proof}
    $S\subseteq\alpha$ and $C\subseteq\alpha$
        by \auto{stationary_in_elim,club_in_elim}.
    Hence $S\inter C\subseteq\alpha$ by \obv{subseteq,inter}.
    Show that for every $D$ such that $\clubIn{D}{\alpha}$
        we have $(S\inter C)\meets D$.
    \begin{subproof}
        Fix $D$.
        Assume $\clubIn{D}{\alpha}$.
        $\clubIn{C\inter D}{\alpha}$
            by \obv{two_club_intersection}.
        Hence $S\meets(C\inter D)$ by \obv{stationary_in}.
        Take $x$ such that $x\in S$ and $x\in C\inter D$
            by \auto{disjoint}.
        Hence $x\in C$ and $x\in D$
            by \obv{inter_elim_left,inter_elim_right}.
        Hence $x\in S\inter C$ by \obv{inter_intro}.
        Follows by \auto{disjoint}.
    \end{subproof}
    Follows by \obv{stationary_in}.
\end{proof}

\begin{proposition}\label{limit_ordinal_is_club_in_itself}
    Let $\alpha$ be a limit ordinal.
    Then $\clubIn{\alpha}{\alpha}$.
\end{proposition}
\begin{proof}
    Show that $\ordinalUnbounded{\alpha}{\alpha}$.
    \begin{subproof}
        $\alpha\subseteq\alpha$ by \obv{subseteq_refl}.
        For every $\beta\in\alpha$ there exists $\gamma\in\alpha$ such that
            $\beta\in\gamma$ by \auto{suc_elem_limitordinal,suc_intro_self}.
        Follows by \obv{ordinal_unbounded}.
    \end{subproof}
    Show that $\ordinalClosed{\alpha}{\alpha}$.
    \begin{subproof}
        $\alpha\subseteq\alpha$ by \obv{subseteq_refl}.
        For every limit ordinal $\gamma\in\alpha$ such that
            $\ordinalUnbounded{\alpha\inter\gamma}{\gamma}$
            we have $\gamma\in\alpha$.
        Follows by \obv{ordinal_closed}.
    \end{subproof}
    Follows by \obv{club_in}.
\end{proof}

\begin{theorem}\label{limit_ordinals_below_is_club}
    Let $\alpha$ be a limit ordinal.
    Suppose $\naturals\in\cofinality{\alpha}$.
    Then $\clubIn{\limitOrdinalsBelow{\alpha}}{\alpha}$.
\end{theorem}
\begin{proof}
    Let $L=\limitOrdinalsBelow{\alpha}$.
    $L\subseteq\alpha$ by \auto{limit_ordinals_below,subseteq}.
    Show that $\ordinalUnbounded{L}{\alpha}$.
    \begin{subproof}
        Show that for every $\beta\in\alpha$ there exists $\gamma\in L$
            such that $\beta\in\gamma$.
        \begin{subproof}
            Fix $\beta\in\alpha$.
            $\clubIn{\alpha}{\alpha}$
                by \obv{limit_ordinal_is_club_in_itself}.
            Take $\gamma\in\alpha\inter\alpha$ such that
                $\beta\in\gamma$ and $\gamma$ is a limit ordinal
                by \obv{two_clubs_have_common_point_above}.
            Hence $\gamma\in\alpha$ by \obv{inter_elim_left}.
            Hence $\gamma\in L$ by \obv{limit_ordinals_below}.
            Follows by \auto.
        \end{subproof}
        Follows by \obv{ordinal_unbounded}.
    \end{subproof}
    Show that $\ordinalClosed{L}{\alpha}$.
    \begin{subproof}
        Show that for every limit ordinal $\gamma\in\alpha$ such that
            $\ordinalUnbounded{L\inter\gamma}{\gamma}$
            we have $\gamma\in L$.
        \begin{subproof}
            Fix $\gamma$.
            Assume $\gamma\in\alpha$.
            Assume $\gamma$ is a limit ordinal.
            Assume $\ordinalUnbounded{L\inter\gamma}{\gamma}$.
            Follows by \obv{limit_ordinals_below}.
        \end{subproof}
        Follows by \obv{ordinal_closed}.
    \end{subproof}
    Follows by \auto{club_in}.
\end{proof}

\begin{proposition}\label{normal_map_reflects_order}
    Suppose $\normalMap{c}{\lambda}{\kappa}$.
    Let $i,j\in\lambda$.
    Suppose $c(i)\in c(j)$.
    Then $i\in j$.
\end{proposition}
\begin{proof}
    $\lambda$ is an ordinal by \auto{normal_map_type}.
    $i$ and $j$ are ordinals by \obv{elem_of_ordinal_is_ordinal}.
    $i\in j$ or $j\in i$ or $i=j$ by \obv{ordinal_elem_connex}.
    Show that $j\notin i$.
    \begin{subproof}
        Suppose $j\in i$.
        Hence $c(j)\in c(i)$
            by \obv{normal_map_increasing_elim}.
        Contradiction by \obv{in_asymmetric}.
    \end{subproof}
    $i\neq j$ by \auto{in_irrefl}.
\end{proof}

\begin{theorem}\label{normal_cofinal_map_range_is_club}
    Let $\kappa$ be a limit ordinal.
    Suppose $\normalCofinalMap{c}{\lambda}{\kappa}$.
    Then $\clubIn{\ran{c}}{\kappa}$.
\end{theorem}
\begin{proof}
    $c\in\funs{\lambda}{\kappa}$ and
        $\cofinalMap{c}{\lambda}{\kappa}$
        by \obv{normal_cofinal_map_type,normal_cofinal_map_is_cofinal}.
    $\lambda$ is an ordinal by \obv{normal_cofinal_map_type}.
    $\normalMap{c}{\lambda}{\kappa}$
        by \obv{normal_cofinal_map_is_normal}.
    Hence $c$ is a function and $\dom{c}=\lambda$
        by \obv{funs_is_function,funs}.
    $\ran{c}\subseteq\kappa$ by \obv{funs_ran}.
    $\ordinalUnbounded{\ran{c}}{\kappa}$
        by \obv{cofinal_map_into_limit_has_unbounded_range}.
    Show that $\ordinalClosed{\ran{c}}{\kappa}$.
    \begin{subproof}
        Show that for every limit ordinal $\gamma\in\kappa$ such that
            $\ordinalUnbounded{(\ran{c})\inter\gamma}{\gamma}$
            we have $\gamma\in\ran{c}$.
        \begin{subproof}
            Fix $\gamma$.
            Assume $\gamma\in\kappa$.
            Assume $\gamma$ is a limit ordinal.
            Assume $\ordinalUnbounded{(\ran{c})\inter\gamma}{\gamma}$.
            Let $A=\{i\in\lambda\mid c(i)\in\gamma\}$.
            $A\subseteq\lambda$ by \obv{subseteq}.
            Show that $A$ is inhabited.
            \begin{subproof}
                $\emptyset\in\gamma$ by \obv{limit_ordinal}.
                Take $y\in(\ran{c})\inter\gamma$ such that
                    $\emptyset\in y$ by \obv{ordinal_unbounded}.
                Hence $y\in\ran{c}$ and $y\in\gamma$
                    by \obv{inter_elim_left,inter_elim_right}.
                Take $i\in\dom{c}$ such that $c(i)=y$
                    by \auto{fun_ran_iff}.
                Hence $i\in\lambda$.
                Hence $c(i)\in\gamma$.
                Hence $i\in A$.
                Follows by \auto.
            \end{subproof}
            Show that for every $i\in A$ there exists $j\in A$ such that
                $i\in j$.
            \begin{subproof}
                Fix $i\in A$.
                Hence $i\in\lambda$ and $c(i)\in\gamma$.
                Take $y\in(\ran{c})\inter\gamma$ such that $c(i)\in y$
                    by \obv{ordinal_unbounded}.
                Hence $y\in\ran{c}$ and $y\in\gamma$
                    by \obv{inter_elim_left,inter_elim_right}.
                Take $j\in\dom{c}$ such that $c(j)=y$
                    by \auto{fun_ran_iff}.
                Hence $j\in\lambda$.
                Hence $c(j)\in\gamma$.
                Hence $j\in A$.
                $c(i)\in c(j)$.
                Hence $i\in j$
                    by \obv{normal_map_reflects_order}.
                Follows by \auto.
            \end{subproof}
            Let $l=\unions{A}$.
            Every element of $A$ is an ordinal
                by \obv{subseteq,elem_of_ordinal_is_ordinal}.
            Hence $l$ is an ordinal by \auto{union_of_ordinals_is_ordinal}.
            Show that $A\subseteq l$.
            \begin{subproof}
                Show that every element of $A$ is an element of $l$.
                \begin{subproof}
                    Fix $i\in A$.
                    Take $j\in A$ such that $i\in j$.
                    Follows by \auto{unions_iff}.
                \end{subproof}
                Follows by \obv{subseteq}.
            \end{subproof}
            Hence $l$ is inhabited by \auto.
            Hence $l\neq\emptyset$ by \auto{nonempty_inhabited}.
            Hence $\emptyset\in l$
                by \obv{ordinal_empty_or_emptyset_elem}.
            Show that for every $r\in l$ we have $\suc{r}\in l$.
            \begin{subproof}
                Fix $r\in l$.
                Take $i\in A$ such that $r\in i$ by \auto{unions_iff}.
                Take $j\in A$ such that $i\in j$.
                Hence $r$ is an ordinal by \obv{elem_of_ordinal_is_ordinal}.
                Hence $\suc{r}\subseteq i$ by \obv{ordinal_suc_subseteq}.
                $\suc{r}$ is an ordinal by \obv{suc_ordinal}.
                Hence $\suc{r}\in i$ or $\suc{r}=i$
                    by \obv{ordinal_subseteq_implies_elem_or_eq}.
                Hence $\suc{r}\in j$ by \auto{ordinal_transitivity}.
                Follows by \auto{unions_iff}.
            \end{subproof}
            Hence $\emptyset\precedes l$.
            Hence $l$ is a limit ordinal by \obv{limit_ordinal}.
            $\suc{\gamma}\in\kappa$ by \obv{suc_elem_limitordinal}.
            $\cofinalIn{\ran{c}}{\kappa}$ by \obv{cofinal_map}.
            Take $y_0\in\ran{c}$ such that $\suc{\gamma}\subseteq y_0$
                by \obv{cofinal_in}.
            Take $i_0\in\dom{c}$ such that $c(i_0)=y_0$
                by \auto{fun_ran_iff}.
            Hence $i_0\in\lambda$.
            $\gamma\in\suc{\gamma}$ by \obv{suc_intro_self}.
            Hence $\gamma\in c(i_0)$ by \auto{subseteq}.
            Show that $A\subseteq i_0$.
            \begin{subproof}
                Show that every element of $A$ is an element of $i_0$.
                \begin{subproof}
                    Fix $i\in A$.
                    Hence $i\in\lambda$ and $c(i)\in\gamma$.
                    $c(i_0)\in\kappa$ by \obv{funs_type_apply}.
                    Hence $c(i_0)$ is an ordinal
                        by \auto{elem_of_ordinal_is_ordinal,limit_ordinal}.
                    Hence $c(i)\in c(i_0)$
                        by \auto{ordinal_transitivity}.
                    Follows by
                        \obv{normal_map_reflects_order}.
                \end{subproof}
                Follows by \obv{subseteq}.
            \end{subproof}
            Show that $l\subseteq i_0$.
            \begin{subproof}
                Show that every element of $l$ is an element of $i_0$.
                \begin{subproof}
                    Fix $r\in l$.
                    Take $i\in A$ such that $r\in i$
                        by \auto{unions_iff}.
                    Hence $i\in i_0$ by \obv{subseteq}.
                    $i_0$ is an ordinal
                        by \obv{elem_of_ordinal_is_ordinal}.
                    Follows by \auto{ordinal_transitivity}.
                \end{subproof}
                Follows by \obv{subseteq}.
            \end{subproof}
            $i_0$ is an ordinal by \obv{elem_of_ordinal_is_ordinal}.
            Hence $l\in i_0$ or $l=i_0$
                by \obv{ordinal_subseteq_implies_elem_or_eq}.
            Hence $l\in\lambda$ by \auto{ordinal_transitivity}.
            Show that $\ordinalUnbounded{A}{l}$.
            \begin{subproof}
                For every $r\in l$ there exists $i\in A$ such that
                    $r\in i$ by \auto{unions_iff}.
                Follows by \obv{ordinal_unbounded}.
            \end{subproof}
            $l\subseteq\lambda$ by \obv{ordinal_elem_implies_subseteq}.
            Let $g=\restrl{c}{l}$.
            $g\in\funs{l}{\kappa}$ by \auto{funs_restrl}.
            Hence $g$ is a function and $\dom{g}=l$
                by \obv{funs_is_function,funs}.
            For every $i\in l$ we have $g(i)=c(i)$
                by \auto{restrl_of_function_apply}.
            Let $U=\unions{\ran{g}}$.
            Show that $U\subseteq\gamma$.
            \begin{subproof}
                Show that every element of $U$ is an element of $\gamma$.
                \begin{subproof}
                    Fix $r\in U$.
                    Take $z\in\ran{g}$ such that $r\in z$
                        by \auto{unions_iff}.
                    Take $i\in\dom{g}$ such that $g(i)=z$
                        by \auto{fun_ran_iff}.
                    Hence $i\in l$.
                    Take $j\in A$ such that $i\in j$
                        by \obv{ordinal_unbounded}.
                    Hence $j\in\lambda$ and $c(j)\in\gamma$.
                    Hence $c(i)\in c(j)$
                        by \obv{normal_cofinal_map_increasing}.
                    $g(i)=c(i)$.
                    Hence $z\in c(j)$.
                    $\gamma$ is an ordinal by \auto{limit_ordinal}.
                    Hence $z\in\gamma$ by \auto{ordinal_transitivity}.
                    Hence $r\in\gamma$ by \auto{ordinal_transitivity}.
                \end{subproof}
                Follows by \obv{subseteq}.
            \end{subproof}
            Show that $\gamma\subseteq U$.
            \begin{subproof}
                Show that every element of $\gamma$ is an element of $U$.
                \begin{subproof}
                    Fix $r\in\gamma$.
                    Take $z\in(\ran{c})\inter\gamma$ such that $r\in z$
                        by \obv{ordinal_unbounded}.
                    Hence $z\in\ran{c}$ and $z\in\gamma$
                        by \obv{inter_elim_left,inter_elim_right}.
                    Take $i\in\dom{c}$ such that $c(i)=z$
                        by \auto{fun_ran_iff}.
                    Hence $i\in\lambda$.
                    Hence $c(i)\in\gamma$.
                    Hence $i\in A$.
                    Hence $i\in l$ by \obv{subseteq}.
                    Hence $g(i)=c(i)$.
                    $g(i)\in\ran{g}$ by \auto{fun_ran_iff}.
                    Hence $z\in\ran{g}$.
                    Follows by \auto{unions_iff}.
                \end{subproof}
                Follows by \obv{subseteq}.
            \end{subproof}
            Hence $U=\gamma$ by \obv{subseteq_antisymmetric}.
            $c(l)=\unions{\ran{\restrl{c}{l}}}$
                by \obv{normal_cofinal_map_continuous}.
            Hence $c(l)=U$.
            Hence $c(l)=\gamma$.
            $c(l)\in\ran{c}$ by \auto{fun_ran_iff}.
        \end{subproof}
        Follows by \obv{ordinal_closed}.
    \end{subproof}
    Follows by \obv{club_in}.
\end{proof}

\begin{theorem}\label{normal_cofinal_map_club_image}
    Let $\kappa$ be a limit ordinal.
    Suppose $\normalCofinalMap{c}{\lambda}{\kappa}$.
    Suppose $\clubIn{D}{\lambda}$.
    Then $\clubIn{\img{c}{D}}{\kappa}$.
\end{theorem}
\begin{proof}
    $\lambda$ is a limit ordinal and $D\subseteq\lambda$ and
        $\ordinalUnbounded{D}{\lambda}$ and
        $\ordinalClosed{D}{\lambda}$
        by \obv{club_in_elim}.
    $c\in\funs{\lambda}{\kappa}$ and
        $\cofinalMap{c}{\lambda}{\kappa}$
        by \obv{normal_cofinal_map_type,normal_cofinal_map_is_cofinal}.
    $\normalMap{c}{\lambda}{\kappa}$
        by \obv{normal_cofinal_map_is_normal}.
    Hence $c$ is a function and $\dom{c}=\lambda$
        by \obv{funs_is_function,funs}.
    $\img{c}{D}\subseteq\ran{c}$ by \obv{img_subseteq_ran}.
    $\ran{c}\subseteq\kappa$ by \obv{funs_ran}.
    Hence $\img{c}{D}\subseteq\kappa$
        by \obv{subseteq_transitive}.
    Show that $\ordinalUnbounded{\img{c}{D}}{\kappa}$.
    \begin{subproof}
        Show that for every $\beta\in\kappa$ there exists
            $z\in\img{c}{D}$ such that $\beta\in z$.
        \begin{subproof}
            Fix $\beta\in\kappa$.
            $\suc{\beta}\in\kappa$ by \obv{suc_elem_limitordinal}.
            $\cofinalIn{\ran{c}}{\kappa}$ by \obv{cofinal_map}.
            Take $y\in\ran{c}$ such that $\suc{\beta}\subseteq y$
                by \obv{cofinal_in}.
            Take $i\in\dom{c}$ such that $c(i)=y$
                by \auto{fun_ran_iff}.
            Hence $i\in\lambda$.
            Take $j\in D$ such that $i\in j$
                by \obv{ordinal_unbounded}.
            Hence $j\in\lambda$ by \obv{subseteq}.
            Hence $c(i)\in c(j)$
                by \obv{normal_cofinal_map_increasing}.
            $c(j)\in\kappa$ by \obv{funs_type_apply}.
            Hence $c(j)$ is an ordinal
                by \auto{elem_of_ordinal_is_ordinal,limit_ordinal}.
            $\beta\in\suc{\beta}$ by \obv{suc_intro_self}.
            Hence $\beta\in c(i)$ by \auto{subseteq}.
            Hence $\beta\in c(j)$ by \auto{ordinal_transitivity}.
            $c(j)\in\img{c}{D}$
                by \auto{img_iff,function_apply_iff}.
            Follows by \auto.
        \end{subproof}
        Follows by \obv{ordinal_unbounded}.
    \end{subproof}
    Show that $\ordinalClosed{\img{c}{D}}{\kappa}$.
    \begin{subproof}
        Show that for every limit ordinal $\gamma\in\kappa$ such that
            $\ordinalUnbounded{(\img{c}{D})\inter\gamma}{\gamma}$
            we have $\gamma\in\img{c}{D}$.
        \begin{subproof}
            Fix $\gamma$.
            Assume $\gamma\in\kappa$.
            Assume $\gamma$ is a limit ordinal.
            Assume $\ordinalUnbounded{(\img{c}{D})\inter\gamma}{\gamma}$.
            $(\img{c}{D})\inter\gamma
                \subseteq(\ran{c})\inter\gamma$
                by \obv{subseteq,inter}.
            $(\ran{c})\inter\gamma\subseteq\gamma$
                by \obv{subseteq,inter}.
            Hence $\ordinalUnbounded{(\ran{c})\inter\gamma}{\gamma}$
                by \obv{ordinal_unbounded_superset}.
            $\clubIn{\ran{c}}{\kappa}$
                by \obv{normal_cofinal_map_range_is_club}.
            Hence $\ordinalClosed{\ran{c}}{\kappa}$
                by \obv{club_in_elim}.
            Hence $\gamma\in\ran{c}$ by \obv{ordinal_closed}.
            Take $l\in\dom{c}$ such that $c(l)=\gamma$
                by \auto{fun_ran_iff}.
            Hence $l\in\lambda$.
            Hence $l$ is an ordinal by \obv{elem_of_ordinal_is_ordinal}.
            Show that $(D\inter l)$ is inhabited.
            \begin{subproof}
                $\emptyset\in\gamma$ by \obv{limit_ordinal}.
                Take $y\in(\img{c}{D})\inter\gamma$ such that
                    $\emptyset\in y$ by \obv{ordinal_unbounded}.
                Hence $y\in\img{c}{D}$ and $y\in\gamma$
                    by \obv{inter_elim_left,inter_elim_right}.
                Take $j\in D$ such that $(j,y)\in c$
                    by \obv{img_iff}.
                Hence $j\in\dom{c}$ by \auto{dom_iff}.
                Hence $j\in\lambda$.
                Hence $c(j)=y$ by \auto{function_apply_iff}.
                Hence $c(j)\in c(l)$.
                Hence $j\in l$
                    by \obv{normal_map_reflects_order}.
                Hence $j\in D\inter l$ by \obv{inter_intro}.
                Follows by \auto.
            \end{subproof}
            Take $j_0$ such that $j_0\in D\inter l$.
            Hence $j_0\in l$ by \obv{inter_elim_right}.
            Hence $l$ is inhabited by \auto.
            Show that $\ordinalUnbounded{D\inter l}{l}$.
            \begin{subproof}
                $D\inter l\subseteq l$ by \obv{subseteq,inter}.
                Show that for every $i\in l$ there exists $j\in D\inter l$
                    such that $i\in j$.
                \begin{subproof}
                    Fix $i\in l$.
                    $l\subseteq\lambda$
                        by \obv{ordinal_elem_implies_subseteq}.
                    Hence $i\in\lambda$ by \obv{subseteq}.
                    Hence $c(i)\in c(l)$
                        by \obv{normal_cofinal_map_increasing}.
                    Hence $c(i)\in\gamma$.
                    Take $y\in(\img{c}{D})\inter\gamma$ such that
                        $c(i)\in y$ by \obv{ordinal_unbounded}.
                    Hence $y\in\img{c}{D}$ and $y\in\gamma$
                        by \obv{inter_elim_left,inter_elim_right}.
                    Take $j\in D$ such that $(j,y)\in c$
                        by \obv{img_iff}.
                    Hence $j\in\dom{c}$ by \auto{dom_iff}.
                    Hence $j\in\lambda$.
                    Hence $c(j)=y$ by \auto{function_apply_iff}.
                    Hence $c(i)\in c(j)$.
                    Hence $i\in j$
                        by \obv{normal_map_reflects_order}.
                    $c(j)\in c(l)$.
                    Hence $j\in l$
                        by \obv{normal_map_reflects_order}.
                    Hence $j\in D\inter l$ by \obv{inter_intro}.
                    Follows by \auto.
                \end{subproof}
                Follows by \obv{ordinal_unbounded}.
            \end{subproof}
            Hence $l$ is a limit ordinal
                by \obv{positive_ordinal_with_unbounded_subset_is_limit}.
            Hence $l\in D$ by \obv{ordinal_closed}.
            $(l,c(l))\in c$ by \obv{function_apply_elim}.
            Hence $c(l)\in\img{c}{D}$ by \auto{img_iff}.
        \end{subproof}
        Follows by \obv{ordinal_closed}.
    \end{subproof}
    Follows by \obv{club_in}.
\end{proof}

\begin{theorem}\label{normal_cofinal_map_stationary_pullback}
    Let $\lambda$ be a limit ordinal.
    Suppose $\normalCofinalMap{c}{\lambda}{\kappa}$.
    Suppose $S$ is stationary in $\kappa$.
    Then $\normalMapPreimage{c}{S}{\lambda}$ is stationary in $\lambda$.
\end{theorem}
\begin{proof}
    Let $T=\normalMapPreimage{c}{S}{\lambda}$.
    $T\subseteq\lambda$ by \auto{normal_map_preimage,subseteq}.
    Show that for every $D$ such that $\clubIn{D}{\lambda}$
        we have $T\meets D$.
    \begin{subproof}
        Fix $D$.
        Assume $\clubIn{D}{\lambda}$.
        $\clubIn{\img{c}{D}}{\kappa}$
            by \obv{normal_cofinal_map_club_image}.
        Hence $S\meets\img{c}{D}$ by \obv{stationary_in}.
        Take $y$ such that $y\in S$ and $y\in\img{c}{D}$
            by \auto{disjoint}.
        Take $i\in D$ such that $(i,y)\in c$ by \obv{img_iff}.
        $c\in\funs{\lambda}{\kappa}$
            by \obv{normal_cofinal_map_type}.
        Hence $c$ is a function by \obv{funs_is_function}.
        Hence $i\in\dom{c}$ and $c(i)=y$
            by \auto{dom_iff,function_apply_iff}.
        Hence $i\in\lambda$ by \auto{funs}.
        Hence $c(i)\in S$.
        Hence $i\in T$ by \obv{normal_map_preimage}.
        Follows by \auto{disjoint}.
    \end{subproof}
    Follows by \auto{stationary_in}.
\end{proof}
